% Page 091

\seite{91}

\margmark{21. März     F}%
rühlingsanfang! Nicht uns noch den Natur=\ob
die auf in der Nation. Die Nächte frieren noch\ob
streng hält, die Tage jedoch haben milder, beinahe\ob
zu warm. Die Augen fangen zu schmelzen\ob
an. Auf den noch fortbestehend einseitigem Wetterlage\ob
führt Heilmond\unsicher. Das erste Grün treibt auf\ob
voll der, alle Bäume und Sträuchen vorstehen\ob
Geburtstag auf ist der Boden aufgestattet. In\ob
unseren Dörfern liegen an einzelnen Stel=\ob
len noch tiefe Sülpfützen. Die ihren Platz wohl\ob
noch für längerer Zeit gepachtet haben. Seit\ob
Weihnachten will die Wasserleitung im\ob
Tiefstand\unsicher von Regen und mehre Quellen. Offen=\ob
sicht noch ist die Leitung zugefroren. Wenn\ob
es auf noch eine schwere Arbeit ist warten und\unsicher,\ob
so wird der Frühling mit einigen Wochen\ob
den Winterschlaf zu überdauern haben! Hoffen wir, daß\ob
das Wetter recht bald schöner wird! —

\margmark{26. März     H}%
eute war Entlassungs=Feier. Zur Entlassung\ob
kamen 3 Knaben u.\ 2\unsicher junge Agitationäre, 2 Kna=\ob
ben und 5 Mädchen.\ob
Die Osterferien dauern vom 27.\ März bis ein=\ob
schließlich 8.\ April.

Die Amtl. Bekanntmachung welche vor dem 6 ten\ob
in dem Volksschulklassen erscheinen muß, hat sie\ob
auch für einmalige Reparatur und auf Kosten,\ob
um dem Versuch als die Hauptstelle so mögli=\ob
chen, da man genügen sind, so hat man daran\ob
denkt, ist die Akte nun richtig, daß 3 na=\ob
mhafte Paare wieder Waffe.

\margmark{9. April     E}%
ingeschrieben in die Schule wurden\ob
jetzt 7 Knaben, 6 Knaben und 1 Mädchen.\ob
Da Lehrer Ringo infolge Krankheit bis 28.\ob
April beurlaubt ist, übernimmt die in der\ob
Zeit Herr Lehrer Burg, Malbergweich, die Ver=\ob
tretung in der Schule Sefferweich.

\margmark{23. April    A}%
m heut. u. am 24.\ wurde Volksunterschrift an\ob
sämtl. Teile. In den vergangenen 3 Wochen\ob
wurden die 5 Tagen wöchentlich je 3 Unterrichts=\ob
stunden durch den Vertreter erhalten, insgesamt

\pend
\begin{verse}
36 Stunden
\end{verse}
\pstart
\pend
